\documentclass[a4paper,12pt]{article}
\usepackage[T2A]{fontenc}
\usepackage[utf8]{inputenc}
\usepackage[russian]{babel}
\usepackage{amsmath,amssymb,mathtools}
\usepackage[default]{lato}
\renewcommand{\familydefault}{\sfdefault}
\usepackage{icomma}
\usepackage[a4paper,top=20mm,bottom=21mm,left=21mm,right=21mm]{geometry}
\usepackage{microtype}
\usepackage{tikz}
\usetikzlibrary{calc,arrows.meta,positioning,shapes.geometric}
\usepackage{tabularx,booktabs}
\usepackage{enumitem}
\usepackage[hidelinks]{hyperref}
\usepackage{parskip}
\usepackage{fancyhdr}
\usepackage{setspace}
\usepackage{xcolor}
\definecolor{accent}{HTML}{167978}
\definecolor{darkaccent}{HTML}{164D50}
\definecolor{textgray}{HTML}{394346}
\definecolor{lightline}{HTML}{A5C4C1}
\color{textgray}
\onehalfspacing
\setlist[enumerate]{leftmargin=*,itemsep=.65em,topsep=.4em}
\setlist[itemize]{leftmargin=*,itemsep=.3em,topsep=.3em}
\setlength{\parindent}{0pt}
\pagestyle{fancy}
\fancyhf{}
\renewcommand{\headrulewidth}{0pt}
\fancyfoot[L]{\footnotesize Лёвин А.~А.\,·\,Анастасия Павлова}
\fancyfoot[R]{\footnotesize\thepage}
\renewcommand{\footrulewidth}{.3pt}
\newcommand{\checkitem}[1]{\noindent\(\square\)\enspace #1\par}
\newcommand{\keyterm}[1]{{\color{darkaccent}\textbf{#1}}}
\newcommand{\smallrule}{\par\vspace{.1em}{\color{lightline}\rule{\linewidth}{.35pt}}\vspace{.1em}}
\tikzset{
 flowstart/.style={ellipse,draw=accent,very thick,align=center,minimum width=30mm,minimum height=9mm,font=\small},
 flowstep/.style={rounded corners=2pt,draw=accent,thick,align=center,text width=105mm,minimum height=12mm,inner sep=6pt,font=\small},
 flowchoice/.style={diamond,aspect=2.5,draw=accent,thick,align=center,text width=37mm,inner sep=4pt,font=\small},
 flowbranch/.style={rounded corners=2pt,draw=accent,thick,align=center,text width=49mm,minimum height=15mm,inner sep=5pt,font=\small},
 flowarrow/.style={-{Stealth[length=2.8mm]},thick,color=darkaccent}
}
\begin{document}
\begin{titlepage}
\thispagestyle{empty}
\begin{center}
\vspace*{42mm}
{\Large\color{accent}\bfseries ЧЕК-ЛИСТ\par}
\vspace{11mm}
{\Huge\color{darkaccent}\bfseries Задачи на смеси,\\ растворы и сплавы\par}
\vspace{7mm}
{\Large Массовая доля, сохранение вещества,\\ таблицы и системы уравнений\par}
\vspace{21mm}
{\large Математика\par}
\vspace{10mm}
{\large 12.10.26\par}
\vfill
{\normalsize Лёвин Артём Александрович\par}
{\normalsize эксклюзивно для Анастасии Павловой\par}
\vspace{20mm}
\begin{tikzpicture}[remember picture,overlay]
 \draw[accent!45,line width=1.05pt] ([yshift=20mm,xshift=22mm]current page.south west)
  .. controls ([yshift=29mm,xshift=68mm]current page.south west) and
             ([yshift=7mm,xshift=-74mm]current page.south east) ..
             ([yshift=20mm,xshift=-22mm]current page.south east);
\end{tikzpicture}
\end{center}
\end{titlepage}
\setcounter{page}{2}

\section*{1. Концентрация: часть и целое}
\checkitem{\keyterm{Масса раствора} складывается из массы растворённого вещества и массы растворителя (обычно воды):}
\[
\boxed{m_{\text{р-ра}}=m_{\text{в-ва}}+m_{\text{воды}}.}
\]
\checkitem{\keyterm{Массовая доля} \(w\) показывает, какую часть массы всего раствора составляет масса вещества. \(p\) --- та же доля, выраженная в процентах:}
\[
\boxed{w=\frac{m_{\text{в-ва}}}{m_{\text{р-ра}}},\qquad
p=100w\%,\qquad m_{\text{в-ва}}=w\,m_{\text{р-ра}}.}
\]
\checkitem{Чтобы перейти от процентов к дроби, делите число процентов на 100: \(15\%\leftrightarrow0{,}15\), \(5\%\leftrightarrow0{,}05\).}
\checkitem{Если хорошо перемешать однородный раствор и взять любую его часть, массовая доля вещества в этой части останется той же.}

\subsection*{Пример 1. Вода и соль}
В 90~г воды растворили 10~г соли. Масса раствора \(90+10=100\)~г, поэтому
\[
w=\frac{10}{100}=0{,}1,\qquad p=10\%.
\]
Если отлить 10~г однородного раствора, в порции будет \(0{,}1\cdot10=1\)~г соли.

\smallrule
\section*{2. Смешивание двух растворов}
\checkitem{При смешивании без потерь складываются \keyterm{массы растворов} и \keyterm{массы вещества}. Сами проценты складывать нельзя.}
\[
\boxed{m_1+m_2=m,\qquad w_1m_1+w_2m_2=wm.}
\]
\checkitem{Удобная таблица: масса смеси, массовая доля, масса нужного вещества. Последнюю колонку получайте умножением двух предыдущих.}

\begin{center}
\renewcommand{\arraystretch}{1.26}
\begin{tabularx}{\linewidth}{@{}Xccc@{}}
\toprule
\textbf{Раствор} & \textbf{Масса, кг} & \textbf{Доля} & \textbf{Вещества, кг}\\
\midrule
Первый & \(4\) & \(0{,}15\) & \(4\cdot0{,}15=0{,}6\)\\
Второй & \(6\) & \(0{,}25\) & \(6\cdot0{,}25=1{,}5\)\\
Смесь & \(10\) & \(x\) & \(10x\)\\
\bottomrule
\end{tabularx}
\end{center}
\subsection*{Пример 2. Два раствора одной соли}
Смешали 4~кг 15\%-го раствора и 6~кг 25\%-го раствора. Тогда
\[
10x=4\cdot0{,}15+6\cdot0{,}25=2{,}1
\quad\Longrightarrow\quad x=0{,}21.
\]
\keyterm{Ответ:} 21\%. Проверка: итоговая концентрация находится между 15\% и 25\%.

\textit{О единицах:} литры измеряют объём, килограммы --- массу. Если в учебной задаче вместо килограммов даны литры, пересчёт требует оговорённого допущения о плотностях и сложении объёмов. В приведённом примере используются массы.

\section*{3. Разбавление и испарение воды}
\checkitem{Чистая вода имеет нулевую долю растворённого вещества: \(w_{\text{воды}}=0\). При добавлении воды масса вещества сохраняется, общая масса увеличивается.}
\checkitem{Если испаряется \keyterm{только вода}, масса растворённого вещества также сохраняется, но общая масса уменьшается.}
\[
\boxed{w_1m_1=w_2m_2\quad\text{(масса вещества сохраняется)}.}
\]
\subsection*{Пример 3. Добавили воду}
В 8~кг 15\%-го раствора добавили 4~кг чистой воды. Получилось 12~кг смеси:
\[
12x=8\cdot0{,}15+4\cdot0=1{,}2
\quad\Longrightarrow\quad x=0{,}1.
\]
\keyterm{Ответ:} 10\%. При разбавлении концентрация уменьшается.

\clearpage
\section*{4. Одинаковые массы и сплавы}
\checkitem{Если смешивают равные массы \(m>0\) двух растворов, можно обозначить их одним символом и сократить:}
\[
2mw=mw_1+mw_2\quad\Longrightarrow\quad
\boxed{w=\frac{w_1+w_2}{2}}.
\]
\checkitem{\keyterm{Сплавы} рассматриваются так же: вместо соли сохраняется масса выбранного металла, например никеля.}
\subsection*{Пример 4. Две равные порции}
Смешали равные массы 15\%-го и 19\%-го растворов. Концентрация смеси равна
\[
\frac{0{,}15+0{,}19}{2}=0{,}17=17\%.
\]
\subsection*{Пример 5. Система для сплавов}
В первом сплаве 20\% никеля, во втором --- 5\%. Из них получили 100~кг сплава с 11\% никеля. Пусть массы исходных сплавов --- \(x\) и \(y\) кг.
\[
\begin{cases}
x+y=100,\\
0{,}20x+0{,}05y=0{,}11\cdot100.
\end{cases}
\quad\Longrightarrow\quad x=40,\;y=60.
\]
Второй сплав тяжелее первого на \(60-40=20\)~кг.

\section*{5. Сушка: сохраняется сухое вещество}
\checkitem{В задачах на высушивание важно выбрать \keyterm{неизменную часть}: массу сухого вещества. Влага уходит, а масса сухого вещества по условию сохраняется.}
\checkitem{Если влажность составляет \(q\%\), доля сухого вещества равна \(1-q/100\).}
\subsection*{Пример 6. Виноград и изюм}
В винограде 90\% влаги, в изюме --- 5\%. Сколько винограда нужно для получения 20~кг изюма? Доля сухого вещества: в винограде \(0{,}10\), в изюме \(0{,}95\).
\[
0{,}10x=0{,}95\cdot20=19
\quad\Longrightarrow\quad x=190.
\]
\keyterm{Ответ:} 190~кг винограда. В обеих частях равенства посчитана именно \keyterm{масса сухого вещества}.

\clearpage
\section*{6. Два эксперимента --- два уравнения}
\checkitem{Если в условии описаны два разных смешивания одних и тех же исходных растворов, запишите \keyterm{по одному уравнению для каждого опыта}. Исходные массы и концентрации при этом общие.}
\subsection*{Пример 7. Вода или ещё один раствор}
Пусть \(x\)~кг 40\%-го раствора и \(y\)~кг 70\%-го раствора сначала смешали с 20~кг воды: получили 41\%-ю смесь. Если заменить воду 20~кг 60\%-го раствора, получится 53\%-я смесь.
\[
\begin{cases}
0{,}4x+0{,}7y=0{,}41(x+y+20),\\
0{,}4x+0{,}7y+0{,}6\cdot20=0{,}53(x+y+20).
\end{cases}
\]
Вычитаем первое уравнение из второго:
\[
12=0{,}12(x+y+20)\ \Longrightarrow\ x+y=80.
\]
Подстановка в первое уравнение даёт \(0{,}4x+0{,}7y=41\). Следовательно, \(x=50\), \(y=30\).

\checkitem{Если неизвестны \keyterm{концентрации} исходных растворов, обозначайте их \(u\) и \(v\). При смешивании \(m\) кг первого и \(n\) кг второго используйте \(mu+nv=(m+n)w\).}
\checkitem{Перед записью ответа проверьте: \(m\ge0\), \(0\le w\le1\); сохранение вещества; соответствие найденной величины вопросу.}

\clearpage
\section*{Алгоритм: как составлять уравнения для смесей}
\vspace{4mm}
\begin{center}
\begin{tikzpicture}[node distance=10mm]
\node[flowstart] (a) {Начало};
\node[flowstep,below=of a] (b) {Выберите сохраняемую часть: соль, кислоту, металл или сухое вещество};
\node[flowstep,below=of b] (c) {Обозначьте неизвестные; переведите проценты в десятичные доли};
\node[flowchoice,below=12mm of c] (d) {В условии два опыта?};
\node[flowbranch,below left=15mm and 8mm of d] (l) {Один опыт: составьте одну таблицу и уравнение};
\node[flowbranch,below right=15mm and 8mm of d] (r) {Два опыта: составьте две таблицы и систему};
\node[flowstep,below=31mm of d] (e) {Решите уравнение или систему. Найдите именно запрошенную величину};
\node[flowstep,below=of e] (f) {Проверьте массы, доли, единицы и сохранение вещества};
\node[flowstart,below=of f] (g) {Ответ};
\draw[flowarrow] (a)--(b);
\draw[flowarrow] (b)--(c);
\draw[flowarrow] (c)--(d);
\draw[flowarrow] (d.south west) -| node[near start,left,font=\footnotesize]{нет} (l.north);
\draw[flowarrow] (d.south east) -| node[near start,right,font=\footnotesize]{да} (r.north);
\draw[flowarrow] (l.south) -- ([xshift=-29mm]e.north);
\draw[flowarrow] (r.south) -- ([xshift=29mm]e.north);
\draw[flowarrow] (e)--(f);
\draw[flowarrow] (f)--(g);
\end{tikzpicture}
\end{center}

\clearpage
\section*{Тренировочный блок \hfill Путь А}
\textit{10 заданий. Для задач со смесями составьте таблицу «масса --- доля --- масса вещества». Записывайте уравнение, вычисления и ответ с единицами.}
\begin{enumerate}
\item В 180~г воды растворили 20~г соли. Найдите массовую долю соли в полученном растворе (в процентах).
\item Смешали 3~кг 10\%-го раствора и 7~кг 30\%-го раствора одного вещества. Найдите концентрацию смеси.
\item К 5~кг 24\%-го раствора соли добавили 3~кг чистой воды. Найдите концентрацию нового раствора.
\item Смешали равные массы 18\%-го и 26\%-го растворов. Сколько процентов составляет концентрация смеси? Объясните, почему общая масса не нужна.
\item В сосуде находилось 12~кг 25\%-го раствора соли. Из него испарилось 3~кг воды, соль осталась в сосуде. Найдите новую концентрацию соли.
\item Первый сплав содержит 20\% никеля, второй --- 50\%. Из них изготовили 30~кг сплава с 32\% никеля. Сколько килограммов каждого сплава использовали? На сколько килограммов их массы различаются?
\item В свежих ягодах 80\% воды, а в высушенных --- 10\%. Сколько килограммов свежих ягод требуется для получения 8~кг высушенных? Считайте, что при сушке уходит только вода.
\item Смешали \(x\)~кг 40\%-го и \(y\)~кг 70\%-го растворов кислоты. Если добавить к ним 10~кг воды, получается 41\%-й раствор; если вместо воды добавить 10~кг 60\%-го раствора той же кислоты, получается 53\%-й раствор. Найдите \(x\) и \(y\).
\item В первом сосуде 6~кг раствора кислоты, во втором --- 12~кг. При смешивании всего содержимого получается 40\%-й раствор. При смешивании \emph{равных масс} исходных растворов получается 35\%-й раствор. Сколько килограммов кислоты было в первом сосуде?
\item В сосуде было 30~кг 20\%-го раствора кислоты. Из него выпарили \(x\)~кг воды, затем добавили \(x\)~кг 60\%-го раствора той же кислоты. Концентрация полученной смеси составила 35\%. Найдите \(x\). Учтите, что масса конечного раствора равна исходной.
\end{enumerate}

\clearpage
\section*{Ответы для самопроверки}
\renewcommand{\arraystretch}{1.65}
\begin{tabularx}{\linewidth}{@{}cX@{}}
\toprule
\textbf{№} & \textbf{Ответ}\\
\midrule
1 & 10\%.\\
2 & 24\%.\\
3 & 15\%.\\
4 & 22\%.\\
5 & \(33\dfrac13\%\) (или \(\dfrac13\) всей массы).\\
6 & 18~кг первого, 12~кг второго; разница 6~кг.\\
7 & 36~кг свежих ягод.\\
8 & \(x=25\)~кг, \(y=15\)~кг.\\
9 & 1{,}2~кг кислоты в первом сосуде.\\
10 & \(x=7{,}5\)~кг.\\
\bottomrule
\end{tabularx}
\vspace{8mm}
\textit{Контрольный вопрос: какая величина сохраняется при каждом изменении смеси?}
\end{document}
